%  LaTeX support: latex@mdpi.com 
%  For support, please attach all files needed for compiling as well as the log file, and specify your operating system, LaTeX version, and LaTeX editor.

%=================================================================
\documentclass[land,article,submit,pdftex,moreauthors]{Definitions/mdpi} 

%=================================================================
%----------
% journal
%----------
% Choose between the following MDPI journals:

% MDPI internal commands
\firstpage{1} 
\makeatletter 
\setcounter{page}{\@firstpage} 
\makeatother
\pubvolume{1}
\issuenum{1}
\articlenumber{0}
\pubyear{2023}
\copyrightyear{2023}
%\externaleditor{Academic Editor: Firstname Lastname}
\datereceived{22 September 2023} 
\daterevised{} 
\dateaccepted{} 
\datepublished{} 
\hreflink{https://doi.org/} % If needed use \linebreak
%\doinum{}

%=================================================================
% Add packages and commands here. The following packages are loaded in our class file: fontenc, inputenc, calc, indentfirst, fancyhdr, graphicx, epstopdf, lastpage, ifthen, lineno, float, amsmath, setspace, enumitem, mathpazo, booktabs, titlesec, etoolbox, tabto, xcolor, soul, multirow, microtype, tikz, totcount, changepage, attrib, upgreek, cleveref, amsthm, hyphenat, natbib, hyperref, footmisc, url, geometry, newfloat, caption

\graphicspath{{figures/}} % path to folder with figures
\usepackage{subcaption}
\usepackage{listings}
\usepackage{xcolor}
\usepackage{gensymb} % for \textdegree
\usepackage{tabularx}
\usepackage{makecell}
\usepackage{multirow}
\usepackage{verbatim}

%=================================================================
% Full title of the paper (Capitalized)
\Title{Monitoring Seasonal Fluctuations in Saline Lakes of Tunisia Using Earth Observation Data Processed by GRASS GIS}
%Monitoring Seasonal Fluctuations in Sabkhas of Tunisia by Multispectral Image Analysis Using GRASS GIS

% MDPI internal command: Title for citation in the left column
\TitleCitation{Monitoring Seasonal Fluctuations in Saline Lakes of Tunisia Using Earth Observation Data Processed by GRASS GIS}

% Author Orchid ID: enter ID or remove command
\newcommand{\orcidauthorA}{0000-0002-5759-1089} % Polina Add \orcidA{} behind the author's name

% Authors, for the paper (add full first names)
\Author{Polina Lemenkova $^{1}$*\orcidA{}}

%\longauthorlist{yes}

% MDPI internal command: Authors, for metadata in PDF
\AuthorNames{Polina Lemenkova}

% MDPI internal command: Authors, for citation in the left column
\AuthorCitation{Lemenkova, P.}
% If this is a Chicago style journal: Lastname, Firstname, Firstname Lastname, and Firstname Lastname.

% Affiliations / Addresses (Add [1] after \address if there is only one affiliation.)
\address{%
$^{1}$ \quad Department of Geoinformatics, \hl{Faculty of Digital and Analytical Sciences}, Universität Salzburg, Schillerstraße 30, A-5020 Salzburg, Austria}

% Contact information of the corresponding author
\corres{Correspondence: polina.lemenkova@plus.ac.at; Tel.: +43-677-6173-2772 (P.L.)}

% Current address and/or shared authorship
%\firstnote{Current address: Affiliation 3.} 

% Abstract (Do not insert blank lines, i.e. \\) 
\abstract{This study documents the changes of the Land Use/Land Cover (LULC) in the region of saline lakes in north Tunisia, Sahara Desert. \textcolor{blue}{Remote} sensing data are valuable data source to monitoring LULC in lacustrine landscapes, because variations in the extent of lakes are visible from space and can be detected on the images. In this study, changes in LULC of the salt pans of Tunisia \textcolor{blue}{are} evaluated using a series of 12 Landsat 8-9 OLI/TIRS. The images were processed with GRASS GIS software. The study area included four salt lakes of north Tunisia in the two regions of the Gulf of Hammamet and Gulf of Gabès: 1) Sebkhet de Sidi el Hani (Sousse Governorate), 2) Sebkha de Moknine (Mahdia Governorate), 3) Sebkhet El Rharra and 4) Sebkhet en Noual (Sfax). \textcolor{blue}{A quantitative estimate of the areal extent analyzed in this study is 182 km x 185 km for each Landsat scene in two study areas: Gulf of Hammamet and Gulf of Gabès.} The images were analysed for \textcolor{blue}{the period} 2017--2023 on months February, April and July for each year. Spatio-temporal changes in LULC and their climate-environmental driving forces were analysed. The results were interpreted and the highest changes detected by accuracy assessment, computing the class separability matrices, evaluating the means and standard deviation for each band, and plotting the reject probability maps. Multi-temporal changes in LULC classes are reported for each image. The results demonstrated changes in salt lakes were determined for winter/spring/summer months as detected changes in water/land/salt/sand/vegetation areas. This paper contributes to the environmental monitoring of Tunisian landscapes and analysis of climate effects on LULC in landscapes of north Africa.}

% Keywords
\keyword{Remote sensing; image processing; Africa; \textcolor{blue}{Sahara; image analysis}; Landsat} 

% The fields PACS, MSC, and JEL may be left empty or commented out if not applicable
\PACS{91.10.Da; 91.10.Jf; 91.10.Sp; 91.10.Xa; 96.25.Vt; 91.10.Fc; 95.40.+s; 95.75.Qr; 95.75.Rs; 42.68.Wt}
\MSC{86A30; 86-08; 86A99; 86A04}
\JEL{Y91; Q20; Q24; Q23; Q3; Q01; R11; O44; O13; Q5; Q51; Q55; N57; C6; C61}

\begin{document}

%----------- section ----------->
\section{Introduction}

\subsection{Background}

Remote sensing (RS) data analysis is a fundamental issue in environmental studies and ecosystem monitoring. Processing RS data is especially essential for multi-temporal analysis of land cover changes \cite{BOUSSETTA2022102564} and application of environmental geomorphology, which are subject to seasonal changes and shrinkage due to climate effects \cite{JonesMillington}. Climate effects can cause desertification \cite{Pontanier,info14040249} and threaten balance of biodiversity \cite{YAICHEACHOUR2023501}. Therefore, monitoring environmental properties of Saharan landscapes by RS data plays an essential role in analysis and decision making on environmental sustainability, which results in many approaches to the effective methods of image processing for environmental tasks \cite{Monget,Lemenkova202227,CAMPOS2018211}. 

The key factors in effective approaches to satellite image processing include the \hl{resolution of data (e.g., 5-30 m/pixel for USGS/NASA Landsat  or high resolution with 1-5 m/pixel for Planet Labs Rapid Eye Satellite) }and the functionality of software \cite{rs14236017}. Another important factor consists in the quality of image, coverage extent, cloudiness and haze. Major image types \hl{are} widely used for effective environmental monitoring and include Landsat \cite{BOUSSEMA2023185,SMIDA2022104643,jmse11040871}, Advanced Very High Resolution Radiometer (AVHRR) \cite{Bryant}, Google Earth \cite{LI2022252}, Moderate Resolution Imaging Spectroradiometer (MODIS) \cite{TOUHAMI2022103804,RHIF2022101596}, and the combinations thereof \cite{REBILLARD1984133}. Examples of very high resolution data include such images as Quickbird \cite{CHOUARI2021443}, Satellite Pour l’Observation de la Terre (SPOT) \hl{with spatial resolution ranging from 2.5 to 5 m in panchromatic mode, 10 m in multispectral mode and 20 m on short wave infrared mode, }\hl{, Sentinel-2 mission with spatial resolution of 10 m in 4 bands, 20 m in 6 bands and 60 m 3 bands, and a swath width of 290 km)} and IKONOS \hl{imagery with spatial resolution of 0.80 m in panchromatic and 3.2 m in multispectral bands}\cite{VELA2008591}. Certainly, although diverse images can be used \hl{as }data sources for environmental monitoring and mapping, the general approach of image analysis relies on using imagery as Landsat or other existing sensor types to extract information through image processing \cite{LiangWang}. 

Selecting a suitable GIS is important for \hl{satellite image processing}, since its functionality is necessary for data \hl{capture and information retrieval from the images}. A comprehensive survey is available on the existing free GIS \cite{LEIDIG201549,CHEN2010253}, recent trends in the development of cartography \cite{WangWu} and new directions and perspectives in GIS \cite{Lietal2022}, while this section covers the most relevant GIS applications for environmental mapping. The efficient solutions for satellite image processing are presented by Earth Resources Data Analysis System (ERDAS) Imagine with applications of environmental monitoring \cite{FALAKI2020363} with many existing cases due to the functionality and availability of diverse modules aimed at raster data processing. The Global Mapper GIS, developed by the United States Geological Survey (USGS), includes a variety of tools including the Light Detection and Ranging (LiDAR) handling, 3D rendering, watershed delineation to process geological and environmental data. Probably the most-well known and widely distributed commercial GIS is the ArcGIS which process well both vector and raster data \cite{AlDjazouli}, supports satellite images processing \cite{SESNIE20082145}, enables hydrological catchment and basin analysis \cite{Nkiaka}, terrain DEM modelling for geospatial data analysis \cite{Kay} and spatial data visualization. 

Nevertheless, in order to reduce the uncertainty and increase the effectiveness of environmental monitoring, a lot of approaches to RS data are proposed as an alternative to processing climate data, some of them using Landsat products \cite{Lemenkova202315} and their integration with Google Earth Engine platform \cite{rs15133257}, or mixed approaches with GIS-based mapping \cite{SAIDI20091683,ANANE201236,Lemenkova202217}, and some of them using hydrological datasets for evaluating physicochemical parameters of groundwater in inflow \cite{BENBRAHIM2021104224}. Other examples include evaluating the content of soil organic carbon stocks in topsoils using bioclimatic data \cite{BAHRI2022e00561} or monitoring agricultural and irrigated lands and their behaviour during drought periods \cite{MOHAMED2023e01766}. 

\hl{The focus of the present study is salinisation issues in Tunisia, Figure} \ref{fig01}. \hl{Thus,} approaches to environmental monitoring \hl{of saline lakes consist in} optimised data modelling \hl{for} ecological \hl{assessment of lacustrine environment in arid landscape. The primary aim is} to highlight correlations between climate and environmental responses \cite{BESSER2021104134} and to evaluate regional hydrogeological processes\hl{. Specific soil types in arid clime of} Sahara desert \hl{affect} the geochemical composition of groundwater that \hl{in turns }influences surface water quality and vegetation patterns in Tunisia \cite{HAMDI2021102974,RAFAATRIGUI2021104121}. At the same time, various landscapes have different mechanisms of ecosystem functionality controlling the rates of vegetation response to droughts and reaction in soil materials and phases of dissolution–precipitation during salinisation \cite{NAVARROTORRE2023105397}. Therefore, \hl{understanding possible causes and environmental implications, most particularly in Tunisia rely on proper data analysis. The aim of data modelling is to derive information regarding }the parameters of these processes, \hl{to evaluate }the effects of seasonal rise in temperature and soil performance as a sediment background in saline lakes of Tunisia during heat periods \hl{which }is of great importance for environmental monitoring of Tunisian sector of Sahara. \hl{With this regards, remote sensing techniques present a useful tools to monitor at large scale the phenomenon of environmental-climate interactions.}

\hl{In this study,} we explore the advanced cartographic approaches to map changes in saline lakes of Tunisia using satellite images to better evaluate the effects from climate and seasonal fluctuations in temperature on salinisation of sebkhas in north Sahara. The success in satellite image processing can be evaluated through classified and mapped difference in land cover types showing \hl{the }environmental characteristics \hl{of soil types} and land \hl{cover }categories which indicate land cover changes \cite{Shoshany}. The traditional methods of \hl{such }evaluation include \hl{fieldwork} and topographic measurements\hl{. The data obtained from such surveys can then be overlayed on} the existing vector layers using GIS. However, such methods have\hl{ drawbacks: }the organisation \hl{of} fieldwork measurements is required for the whole process of land monitoring, surveying and topographic mapping \cite{JENDOUBI2020100507}. This might be hard and costly \hl{process }to implement in \hl{such} remote areas \hl{as} Sahara desert. Besides, \hl{the }traditional geographic fieldwork measurements provide information of land categories as an discrete values measured by  \hl{technical} surveyors without \hl{sufficient }data on complex patterns of the landscapes. 

An alternative approach is presented by \hl{the remote sensing (RS)} data \hl{processing} which can be used as descriptors of land cover types and partial replacement of the traditional topographic fieldwork methods. \hl{The }use of the RS integral assessment of land cover patterns are possible using satellite images \hl{covering} large extent of the \hl{study} area \cite{ELGHOUL2023106733}. \hl{The} satellite images provide direct information on land categories due to the relationship between spectral reflectance and the environmental properties of land categories \hl{which can be derived from the} spaceborne images \hl{and interpreted using GIS tools}\cite{Allouche,Ruiz,Ray}. As a general rule, vegetation and water areas can be identified using a combination of Red/NIR channels since the reflectance in NIR \hl{increases} for vegetation along with the absorbed signal in Red bands\hl{. This enables} discrimination between water and vegetation areas \hl{during} image analysis. Such RS properties are widely used for calculation of vegetation indices. 

\hl{Generally speaking, the RS} approach is based on measuring \hl{spectral} reflectance of the pixels on the satellite image\hl{. }Spectral reflectance of \hl{pixels} in various wavebands of the multispectral \hl{image }indicates spectral characteristics of different objects on the Earth \hl{and thus represents diverse land cover types}. This fundamental characteristics enables to use \hl{spaceborne }images for detecting land cover types. \hl{Moreover, compensating and correcting the image from atmospheric and air-water effects enable }image interpretation \hl{for thematic mapping}. Therefore, evaluating land cover \hl{changes using} satellite images\hl{ can} be performed using times series approach\hl{. This solution is based on the} interpretation of several satellite images taken on the same target area in different seasons of year or a sequence of several years. 

\hl{With this regard, RS method} emerge as a fundamental technique of Earth Observation in a wide variety of geospatial and environmental applications. \hl{The case studies include }landscape mosaics \cite{BIARNES2021103281}, monitoring sedimentation processes \cite{TOWNSHEND1989177}, dynamics of sandy dunes in the Sahara \cite{Afrasinei}, geomorphological modelling \cite{Desprats} and monitoring coastal regions \cite{AMROUNI201973,Jaquet}. Such generic approach of \hl{RS and advantages of the use of spaceborne data have} been noted earlier \cite{RAMPHERI2023103359}\hl{. These} include \hl{high} sensitivity of satellite images which enables to detect small-size patches of landscapes, regularity of survey by various space sensors and flexibility of data coverage\hl{. The latter enables images to be} finely adjusted for practically any study area using search criteria in \hl{data} selection and threshold parameters \hl{such as }cloudiness, \hl{or} image source. 

When applied to a series of images, processing RS data can effectively discover the environmental properties of  landscapes \cite{Toujani,Quarmby,Gelebo}. Moreover, RS data processing and satellite image analysis \hl{enable }to quantify \hl{landscape} patches within the target \hl{area} through the interpretation of classes \hl{that indicate} land categories and vegetation types \cite{Benedetti,Godinho}. \hl{Furthermore, this approach enables to} determine climate effects in lacustrine or coastal environments \cite{Souissi}. Finally, other advantages of the RS data processing \hl{consist in} a resource- and time-effective \hl{approach} indicating \hl{the }dynamics in landscapes through evaluation of changes \hl{over} time\hl{.} 

Inspired by the existing examples of the RS data analysis in environmental monitoring \cite{f9020059,Lemenkova202229,Fathalli,Abbas,Henchiri}, this study integrated the advanced methods of the GRASS GIS with applied techniques of its image processing modules for assessment and mapping of changes in landscapes of saline lakes in northern Tunisia, Figure \ref{fig01}. The \hl{objective} of the current paper is to provide \hl{the} advanced RS and cartographic methods \hl{to evaluate} the response \hl{of land cover types around salt lakes }to seasonal fluctuations in temperature \hl{and} the \hl{increase} of salinisation. Importantly, the study presented here is not intended to suggest that RS data replace the hydrological datasets. In contrast, these images can supplement the fieldwork survey in the sabkhas and environmental tests on the saline soil of Tunisia. 

\begin{figure}[H]
	\centering
		\includegraphics[width=9.5 cm]{Fig_01a.jpg}
\caption{Two segments of the study area on a topographic map of Tunisia: the Gulf of Hammamet (northern part) and the Gulf of Gabès (southern part). \hl{Yellow dots indicate the locations of the annotated major cities. Two specific locations studied in this research are contoured by purple-coloured rotated square (Gulf of Hammamet) and magenta-coloured rotated square (Gulf of Gab\'es). Mapping software: Generic Mapping Tools (GMT).}}\label{fig01}
\end{figure}

In this paper, the integration of image processing workflow with cartographic and geologic maps provides \hl{the} effective \hl{approach} by which the analysis of environmental changes \hl{can} be performed \hl{with regard to inter-year and multi-decades meteorological variations}. The goal of using the RS data is to qualitatively evaluate and \hl{visualise} changes in land cover types\hl{. The reasons for these changes are }caused by seasonal fluctuations in temperature during winter/spring/summer months and in the process of evaporation. In this way, \hl{RS data} processing \hl{presents} the effective means by which the time series of the multispectral images \hl{can} be analysed for landscape monitoring. Data obtained from \hl{clustering} of the \hl{Landsat OLI/TIRS }satellite images \hl{were} used for modelling land cover types using classification. 

The cartographic experiments \hl{were} performed on two \hl{study} areas in north Tunisia: \hl{1)} the coastal regions of the Gulf of Gabès (salt lake Sebkhet en Noual) and \hl{2) }the Gulf of Hammamet (salt lakes Sebkhet de Sidi el Hani, Sebkha de Moknine and Sebkhet El Rharra). Using a comparison of several images classified for various seasons and time periods, this study \hl{evaluated} the properties of the landscapes using maximum-likelihood discriminant analysis classifier of GRASS GIS. The advanced methods of evaluation of landscape properties are built on the growing demand for the time- and cost-effective approaches to environmental monitoring in changing climate of Tunisia. \hl{This contribute to }saving agricultural resources \hl{and} taking decisions of farmers and environmental planners. \hl{In turn, the effective landscape management enables} to make use of the renewable resources in the \hl{regions of salt lakes considering} regional and local \hl{environmental} settings of \hl{Tunisia.}

\subsection{Study Area}

This study focuses on the salt lakes located in the north Sahara, Tunisia, Figure \ref{fig01}. \hl{Saline lakes in Tunisia are important components of the environment playing a crucial role as biodiversity maintenance, habitats for rare species in nearby oases, rare sources of water, flood and fishery, regulators of temperature, hydrology, and carbon fixation. }Salt lakes, or sebkhas, of Tunisia present a series of saline depressions that are characterised by variability of water: during the winter, sebkhas receive rare water inflow from occasional rains while during summer months the lakes dry out almost completely and become the saline crystallised surfaces \cite{berthon1970saumures,m2001valorisation,Gautier}. 

\hl{The physiographic data on Tunisia present moderate hilly terrain with the slope geomorphology ranging from zero m in the coastal areas at water level to the maximum at 3781 m and mean at 432 m, according to GEBCO DEM grid. }The geographical location of these sebkhas is related to the environmental historical development of north Sahara with related aeolian and fluvial processes \cite{ALI20231,MWIRICHIA2023269}. \hl{Located in arid environment of north Africa, saline lakes and surrounding landscapes experience changes in volume and extent due to the extreme temperatures: during summer heat period, lakes dry out almost entirely and become a crust of salt and minerals due to the evaporated moisture. }Gradual deflation and geomorphic erosion lead to the formation of low individual basins and pans where current salt lakes are located \hl{on specific soil types typical for arid environment }\cite{Blackwelder140}. \hl{Soil classification in Tunisia include the following types of soil according to the the French system: podzols, vertisol, red Mediterranean, calcium-magnesium, brown, saline and hydromorphic very acidic soil, and poorly evolved soils. Of these, the dominant type is the calcium-magnesium soils due to the influence of arid environment and Sahara desert.}

The sebhkas are mostly located in topographic depressions which remain as geomorphic landforms from the earlier palaeolacustrine basins, palaeodrainages, interdunes, and coastal plains \cite{GOUDIE19951}. The formation of these depressions is strongly connected to the tectonic movements in Quaternary which shaped present morphology of sebkhas \cite{Abbes,Louhaichi,Ballais}. The lagoon facies of the selected lakes evident the stratigraphic distribution of gypsum clays and laminated gypsum lithofacies from the Neogene period (Mio--Pliocene and Quaternary evaporites) and Cretaceous sediments \cite{BARBIERI2006726,SWEZEY1996335}. This supports the hypothesis of earlier existence of the Sahara Sea with endorheic basins that nowadays become a series of depressions of the saline lakes in north Africa \cite{Coque}. 

The effects of climate and seasonal meteorological fluctuations on soil and vegetation setting of the Saharan landscapes are explained by changes in temperature which triggers processes of evaporation, i.e., the transform of liquid water in lakes and ephemeral river flows into gaseous vapour \cite{HAMMI2003215}. \hl{Although data on rainfall vary significantly by regions of Tunisia, the average annual rainfall estimates at 158 mm/yr at the country level, with less than 100 mm/yr close to Sahara desert and over 700 mm/yr in the coastal areas. }Natural factors are augmented by the anthropogenic activities that increased the intensity of land cover changes \cite{RODRIGUEZCABALLERO2017146}. \hl{The combination of natural climate and anthropogenic factors results in diversified land cover types classification in Tunisia. Thus, major land cover types include rangelands, riparian vegetation, wastelands, water areas, wetlands and lacustrine areas, artificial and urban areas, cultivated lands and plantations. }High evaporation during the summer months results in a total desiccation of these ephemeral lakes and lead to a formation of the thick layer of minerals and salt on the bottom of the lake basins \cite{BRYANT1994269}. 

Lakes in Tunisia are important components of the dryland environment of Sahara \hl{and strongly depend on the temperature fluctuations. Annual temperature averages in Tunisia reaches its peak between June to September when it exceeds 40\degree C  in southern regions of the country and has low values at 11.5\degree C during the coldest period in January. Hence, salt lakes} are rare sources of water, flood and fishery, they serve as regulators of temperature, hydrology, and they support carbon fixation for environmental sustainability \cite{NIGAM2020201}. Moreover, they play a crucial role for maintaining biodiversity \cite{Neji}, they serve as habitats for rare African species and birds in nearby oases \cite{Baduel,RADDADI2022104806,Kamel} and habitats for halophilic species \cite{NEIFAR20191802}. Such phenomena were revealed in studies on environmental monitoring which evaluated landscape changes in Tunisia. For instance, the reaction of landscapes and vegetation responses to droughts were modelled with regard to the temperature fluctuations and climatic effects in Sahara \cite{HANAFI2008557,PAUSATA2020235,Lemenkova202313}. At the same time, the increase of temperature and reduced precipitation observed during heat of summer months lead to lack of surface water which is the primary factor of desertification and aridification of landscapes in Sahara \cite{OSWALD2023369,DODORICO2013326}. 

%----------- section ----------->
\section{Materials and Methods}

%----------- subsection ----------->
\subsection{Data}

The materials include the 12 Landsat 8-9 OLI/TIRS images collected from the USGS Earth Explorer repository.The USGS Identifier (ID) of the scenes of the satellite images are summarised in Table \ref{tab01}. Within the dataset, 6 images cover the region of the Gulf of Hammamet (Figure \ref{fig03}), and 6 images cover the area of the Gulf of Gabès (Figure \ref{fig04}).

\begin{table}[H]
\caption{\hl{Scene IDs} of the 12 Landsat 8-9 OLI/TIRS images in the EarthExplorer USGS repository}.
\label{tab01}
%\begin{adjustwidth}{-\extralength}{-2.5cm}
%\newcolumntype{L}{>{\centering\scriptsize\arraybackslash}l}
%\setlength\extrarowheight{-0.8cm}
\footnotesize
%\begin{tabularx}{\fulllength}{p{1.5cm}|p{6.5cm}|p{6.5cm}|p{1.9cm}}
\begin{tabularx}{7cm}{l|l}

			\toprule
			\textbf{Date} & \textbf{Scene ID} \\
			\midrule
			\multicolumn{2}{c}{\textbf{Gulf of Hammamet region}}\\
			\midrule
			2023/07/15 & LC91910352023196LGN00 \\	
			2023/04/10 & LC91910352023100LGN00 \\	
			2023/02/21 & LC91910352023052LGN01 \\	
			2017/07/06 & LC81910352017187LGN00 \\	
			2017/04/01 & LC81910352017091LGN00 \\	
			2017/02/28 & LC81910352017059LGN00 \\	
			\midrule
			\multicolumn{2}{c}{\textbf{Gulf of Gabès region}}\\
			\midrule
			2023/07/15 & LC91910362023196LGN00 \\	
			2023/04/10 & LC91910362023100LGN00 \\	
			2023/02/21 & LC91910362023052LGN01 \\	
			2017/07/22 & LC81910362017203LGN00 \\
			2017/04/01 & LC81910362017091LGN00 \\
			2017/02/28 & LC81910362017059LGN00 \\
			\bottomrule
		\end{tabularx}
%\end{adjustwidth}
\end{table} 

\hl{The two segments of study area were chosen, to illustrate the difference in the northern and southern environmental setting of the regions surrounding salt lakes in different geographic regions: northern segment was selected for the Gulf of Hammamet, and the southern segment was selected for the Gulf of Gab\'es. Both study areas are reasonably representative for the illustration of the Tunisian geographic setting. The two segments cover the area where potential environmental and climate impacts are most likely from the effects of Sahara (in southern segment) and milder climate of coastal area (in the northern segment). Northern segment is also the region with high intensity of agricultural uses which contrasts with the southern area which is located closer to Sahara desert. The comparative analysis of two diverse segments of Tunisia represents the environmental processes related to coastal zone in the north and arid climate of Sahara in the south.}

All images are selected on the three key control periods: winter with the highest water level (February), spring with transitional period and rains (April) and summer with the peak heat (July). The metadata in the images are listed in the Tables in \hl{submitted auxiliary data}. It is widely accepted that cloudiness and haze has a key impact on the image properties. Therefore, in this study we used the 10\% of cloud coverage for all images. Tested images included spectral bands in visible, Short-wave infrared (SWIR)-1, SWIR2 and Near-infrared (NIR) channels obtained from the OLI sensor. The topographic maps showing two study areas within the country and their enlarged fragments are performed with the Generic Mapping Tools (GMT) software, version  6.1.1 \cite{Wessel} using the existing scripting techniques reported earlier \cite{Lemenkova202212,Lemenkova202206}. 

%----------- subsection ----------->
\subsection{Methods}

The general scheme of the data and methods applied in this study is presented schematically in a flowchart in Figure \ref{fig02}. 

\begin{figure}[H]
\hspace{50pt}\makebox[0.90\linewidth][r]{%
	\centering
		\includegraphics[width=15.0 cm]{Fig_02.jpg}
}
\caption{Flowchart summarising the workflow. Software: R. Diagram source: author.}\label{fig02}
\end{figure}

The Landsat 8-9 OLI/TIRS satellite images were processed with the Geographic Resources Analysis Support System (GRASS) Geographic Information System (GIS) software. Compared to the existing GIS\hl{, }GRASS GIS presents \hl{more} advanced method of geographic visualization \hl{through the use of} scripts \hl{which }optimises \hl{computer} memory and computation resources\hl{. Besides}, the increase of the open spatial data arises a question of the effective tools for their processing. Since GIS have restricted functionality, modern approach \hl{of} GRASS GIS \hl{presents} a powerful alternative through scripting, which \hl{is} a principally new paradigm for cartographic data processing and satellite image analysis. Motivated by \hl{such }powerful functionality of GRASS GIS and its modules for image processing, this study presents a case of \hl{the }processing remote sensing data by \hl{the GRASS GIS }scripts. 

\hl{Firstly,} the images were processed by scripts using clustering using 'i.cluster' module and continued with classification by the 'i.maxlik' module. The scenes were analysed for 10 recent years from 2013 to 2023 using the maximum-likelihood discriminant analysis classifier. \hl{The classification approach is based on the automatic data clustering which performs partitioning the image at pixel level. The pixels are equally distributed throughout space are grouped into clusters based on the similarity of their spectral reflectance values. Using such approach, spatio-temporal} changes \hl{in} water level and salinisation \hl{were analysed}  \hl{over the }target \hl{area of }salt lakes of Tunisia\hl{. }The detection of the saline lakes is possible using spectral properties of the images that enable to distinct highly reflective cyan-coloured areas of the lakes discriminated against the beige-coloured sands and bare land. \hl{The images were uploaded as RGB triplets with varying intensities of components of the selected colors which enabled to visualise these aspects. }Thus, spectral properties of the RS data were evaluated by various GRASS GIS modules and scripting techniques with regards to the water/land/salt discrimination, and demonstrated robust approach to image analysis. 

\subsubsection{Data preprocessing}
\hl{Firstly}, the images were preprocessed by the 'i.landsat.toar' module of GRASS GIS using properties of \hl{the }image and information on sun elevation level to evaluate the reflectance in pixels and calculation of \hl{the }top-of-atmosphere radiance and temperature for all the Landsat OLI scenes, Figure \ref{fig03}. The snippets of GRASS GIS code are concatenated into longer scripts which perform the complete task of mapping. Such\hl{ }partition of the workflow enables the machine to perform selected tasks of image processing in a distributed way. This decreases a possibility of errors and misclassification of the pixels and supports processing \hl{of the }large datasets due to a significant level of automation. 

Hence, the images were corrected and added into GRASS environment using \hl{the }'r.import' module. After the import procedure, the images with all the multispectral bands were grouped in the working directory using 'i.group' module, which is used to select necessary bands, and their data were checked using 'g.list' module. The 'i.group' module of the GRASS GIS was used for processing the structures of the spatial data and grouping them into classes using the k-means algorithm. \hl{The new contributions of such scripting algorithms of GRASS GIS to image processing are as follows: }
\begin{itemize}
	\item \hl{inherent modules for image processing, analysis, spatial modelling, and visualization, graphical plotting and maps production; }
	\item \hl{streamlining and automating a series of duplicated tasks through scripts for iterated processing of different images as a time series of Landsat 8-9 OLI/TIRS data; }
	\item \hl{diversified process of segmented data handling through programming approach that enables to break down the workflow into a set of individual commands and catch and clean data errors before mapping.}
\end{itemize}
\hl{ In this way, the advantages of the automation achieved through the implemented algorithms of GRASS GIS enable to save time when processing a series of Landsat images, reduce data mistake and upscale or downscale the tasks to various regions, accordingly.}

\begin{figure}[H]
\hspace{50pt}\makebox[0.90\linewidth][r]{%
	\begin{subfigure}[b]{.40\textwidth}
		\centering
			\includegraphics[width=0.87\textwidth]{Fig_03a.jpg}
		\caption{28 February 2017}
	\end{subfigure}%
	\begin{subfigure}[b]{.40\textwidth}
		\centering
		\includegraphics[width=0.87\textwidth]{Fig_03b.jpg}
		\caption{1 April 2017}
	\end{subfigure}%
	\begin{subfigure}[b]{.40\textwidth}
		\centering
		\includegraphics[width=0.87\textwidth]{Fig_03c.jpg}
		\caption{6 July 2017}
	\end{subfigure}%
}\\
\vfill \vspace{1mm}
\hspace{50pt}\makebox[0.90\linewidth][r]{%
	\begin{subfigure}[b]{.40\textwidth}
		\centering
			\includegraphics[width=0.87\textwidth]{Fig_03d.jpg}
		\caption{21 February 2023}
	\end{subfigure}%
	\begin{subfigure}[b]{.40\textwidth}
		\centering
		\includegraphics[width=0.87\textwidth]{Fig_03e.jpg}
		\caption{10 April 2023}
	\end{subfigure}%
	\begin{subfigure}[b]{.40\textwidth}
		\centering
		\includegraphics[width=0.87\textwidth]{Fig_03f.jpg}
		\caption{15 July 2023}
	\end{subfigure}%
}
\caption{Landsat images on the Gulf of Hammamet region (northern segment) showing the increase of salt (bright cyan colour) from winter to summer months in sabkhas Sebkhet de Sidi el Hani (Sousse Governorate), Sebkha de Moknine (Mahdia Governorate) and Sebkhet El Rharra (Sfax Governorate): \textbf{(a)}: 28 February 2017; \textbf{(b)}: 1 April 2017; \textbf{(c)}: 6 July 2017; \textbf{(d)}: 21 February 2023; \textbf{(e)}: 10 April 2023; \textbf{(f)}: 15 July 2023. \hl{The images are shown in natural colour composite of Landsat 8 OLI/TIRS based on a band combination of Red (4), Green (3), and Blue (2). Such composite shows healthy vegetation is green, salt lakes in cyan, while soils and bare land in brown.}}\label{fig03}
\end{figure}

The metadata of the images were checked using relevant modules of the GRASS GIS. In particular, the techniques of image processing in GRASS GIS include the auxiliary functions of 'g.list rast' and 'r.info' acting on imported raster bands and showing the identified parameters of images as attributes and a given set of identifiers of \hl{the }technical specifications of the multispectral images: resolution, minimal and maximal values, coordinate system. After \hl{the }image preprocessing, the 'i.group' module was used for grouping the bands before the classification of the Landsat \hl{8-9 OLI-TIRS }images. This approach partitions the image using the embedded algorithm of discriminating the pixels on the images that represent various land cover types in target area using their spectral reflectances. Afterwards, they were clustered using the k-means algorithm embedded in the GRASS GIS which was implemented by the 'i.cluster' module. \hl{Such sequential use of modules is beneficial for creating an automated workflow of image processing, whereas straightforward traditional method would take more time and efforts which is a fail in automation aspect. }

\begin{figure}[H]
\hspace{50pt}\makebox[0.90\linewidth][r]{%
	\begin{subfigure}[b]{.40\textwidth}
		\centering
			\includegraphics[width=0.87\textwidth]{Fig_04a.jpg}
		\caption{28 February 2017}
	\end{subfigure}%
	\begin{subfigure}[b]{.40\textwidth}
		\centering
		\includegraphics[width=0.87\textwidth]{Fig_04b.jpg}
		\caption{1 April 2017}
	\end{subfigure}%
	\begin{subfigure}[b]{.40\textwidth}
		\centering
		\includegraphics[width=0.87\textwidth]{Fig_04c.jpg}
		\caption{22 July 2017}
	\end{subfigure}%
}\\
\vfill \vspace{1mm}
\hspace{50pt}\makebox[0.90\linewidth][r]{%
	\begin{subfigure}[b]{.40\textwidth}
		\centering
			\includegraphics[width=0.87\textwidth]{Fig_04d.jpg}
		\caption{21 February 2023}
	\end{subfigure}%
	\begin{subfigure}[b]{.40\textwidth}
		\centering
		\includegraphics[width=0.87\textwidth]{Fig_04e.jpg}
		\caption{10 April 2023}
	\end{subfigure}%
	\begin{subfigure}[b]{.40\textwidth}
		\centering
		\includegraphics[width=0.87\textwidth]{Fig_04f.jpg}
		\caption{15 July 2023}
	\end{subfigure}%
}
\caption{Landsat images on the Gulf of Gabès region (southern segment) showing the increase of salt (bright cyan colour) from winter to summer months in sabkhas Sebkhet en Noual (Sfax Governorate): \textbf{(a)}: 28 February 2017; \textbf{(b)}: 1 April 2017; \textbf{(c)}: 6 July 2017; \textbf{(d)}: 21 February 2023; \textbf{(e)}: 10 April 2023; \textbf{(f)}: 15 July 2023.}\label{fig04}
\end{figure}

%----------- subsection ----------->
\subsubsection{Data clustering}

During clustering, \hl{the} 'signature file' was created \hl{and} used in the next step to sort pixels into \hl{land} cover classes using \hl{the }'i.cluster' module of the GRASS GIS which creates cluster and covariance matrices. \hl{This classification algorithm partitions the image into clusters and assignment of pixels into each land cover categories. }All \hl{these }classes contained pixels with approximately the same level of spectral reflectance, which \hl{were} automatically evaluated by the \hl{machine-based} analysis of the GRASS GIS\hl{. Such technique differentiates the} number of pixels within each class and thus \hl{assigns the} extent of \hl{the }landscape categories \hl{in each case This enables us} to evaluate the \hl{landscape }patches and their distribution within the scenes under the effects of various seasons on development of salinisation in sabhkas. Besides, the images located in the Gulf of Gabès contained more sandy desert areas compared to the northern segment of the study areas in the Gulf of Hammamet which allowed to compare the effects from geographic distribution of the saline lakes with higher coverage of mixed vegetation in the more temperate climate of north compared to the southern region which is more impacted by the arid climate of Sahara.

%----------- subsection ----------->
\subsubsection{Image classification}

Theoretically, the essential  principle of the land cover type classification is based on the fundamental properties of RS which consists in measurement of spectral reflectance as a reflection of light from the Earth's surface, as can be visible on the scenes, Figure \ref{fig04}. Spectral reflectance signatures of various land cover types (land, water, salt soils, vegetation, urban areas, etc) differs significantly. This is because of the different emitted radiation from these these objects which varies significantly. The Landsat 8-9 OLI/TIRS sensor continuously measures spectral reflectance of the objects on the Earth which are identified as various land cover types. Hence, the performed measurements indicate various different bodies and objects on the Earth using their spectral reflectance visible on the spaceborne satellite images in selected band combinations. Consequently, since measured land cover types varies, it enables to perform classification of the land cover types using RS data processing, following by cartographic visualization and interpretation. 

Practically, the information on spectral reflectance is used for classification procedure. \hl{In this study, the method included a machine learning techniques which is based on a fully automated approach that is independent from human-based sampling. In this case, the partition of the image was relied on computer vision algorithms which determined the sampling categories and assigned pixels to classes automatically using programmed 'i.cluster' module of GRASS GIS. This approach differs significantly from the traditional supervised classification is based on training samples. }Thus, after clustering, the images were processed using special module of the GRASS GIS which performs unsupervised automatic classification of the satellite images. The advanced methods of GRASS GIS present effective instruments for classification -- the 'i.maxlik', developed by M. Shapiro and T. Wen, and supported with additional module by M. Nartiss. The use of the 'i.maxlik' only required to specify the parameters of signature file and related information on groups and subgroups of spectral bands during data partition, and to assign pixels into each land cover class according to the records of the spectral reflectance in binary format. The spectral reflectance of wavelengths which develops proportionally with internal gain in strength in brightness of land cover types identified from space by the sensor. The technical characteristics of this module include two steps: the preceding clustering and following classification. The module is designed to be used for the multispectral images, such as Landsat, using spectral signatures of the pixels. 

The classification procedure was performed with different images (6 scenes for the study area-1 in the Gulf of Gabès, and 6 images for the study area-2 for the Gulf of Hammamet). The aim of the classification to identify cells in the classified image and to group then into classes according to the similarity of spectral reflectance in each given group (class) through estimated spectral reflectance value during automatic image analysis performed by the GRASS GIS algorithm. As for comparison of several images taken on various seasons and various years with 6 year interval, the main idea was to compare the effects of climate, the impact from the Sahara and seasonal meteorological fluctuations on the development of salinisation in sabkhas of Tunisia using comparative analysis. Therefore, the experiments with images were implemented with different calendar seasons and in different regions (central/southern) of Tunisia. 

%----------- subsection ----------->

 \subsubsection{Accuracy analysis}

The images were evaluated for accuracy and pixels with high level of rejection were filtered out using 'r.mapcalc' module of GRASS GIS. The rejection probability values, i.e., the pixel classification confidence levels were computed and visualised and a series of maps. The error matrix and  convergence level of classified pixels were measured using 'i.cluster' module that includes accuracy assessment of classification results in tabular format. After that, the classification was performed by the and the results of the clustering were exported as a series of matrices in CSV format for each images, and the results are reported in the tables in the \hl{auxiliary materials supporting }this study. The classified satellite scenes were verified and class separability matrix was computed for each of the 12 images. The accuracy of the classification was measured and summarised in the tables with a convergence level of classified pixels above 98\% of points stable. Such approach to analyse the accuracy of image processing provided data to evaluate the obtained values of land cover classes at each image of the processed time series.
 
 %----------- subsection ----------->
 \subsubsection{Mapping}

The environmental properties of the landscapes were evaluated with cartographic methods after image processing using visualised maps. The aim is to analyse the seasonal effects of climate (increase of temperature and evaporation) on the behaviour of salt lakes. The complete cartographic workflow used in this study included a sequence of the following modules of GRASS GIS:

\begin{itemize}
	\item 'r.import' using for data import via the embedded GDAL library; the extent and resolution were adjusted to the region of the Landsat scenes;
	\item 'g.list' for listing imported raster data type by the search pattern of GRASS GIS; 
	\item 'g.region' for defining the computational extent of the region to match the scenes; 
	\item 'i.cluster' for clustering of the raster Landsat images using k-means algorithm; 
	\item 'i.maxlik' for unsupervised image classification using maximum likelihood method; 
	\item 'd.mon' for display visualisation and cartographic plotting; 
	\item 'r.colors' for adjusting the colour palette according to the land cover types; 
	\item 'd.rast' for mapping the image in the active graphics frame; 
	\item 'd.legend' for adding the colour legend with textual notations; 
	\item 'd.out.file' for data export to the conventional formats (JPG). 
\end{itemize}

These GRASS GIS modules have been incorporated to perform various steps of cartographic tasks and image processing workflow by using the GRASS GIS console. The land cover types were identified continuously across the images using a sequence of k-means clustering and maximum likelihood classification approaches, and converted to the maps which indicate variations in land cover types over Tunisia during a period of 2017 to 2023 for two image datasets covering the regions on the Gulf of Hammamet and the Gulf of Gabès. The data obtained from the GRASS GIS based image processing workflow were evaluated for environmental analysis.

%----------- section ----------->
\section{Results and Discussion}

The results of the image processing approach of the GRASS GIS demonstrated in this paper show that the summer period speeds up soil salinisation and mineralisation in sabkhas of Tunisia, as demonstrated through mapped series of images on several seasons with increased summer heat, and vice versa. Then, using the information received from spectral reflectance of the pixels on the images and heat development in Sahara during summer months we were able to find the highest degree of salinisation for each lake in various years. Figure \ref{fig05} demonstrates the results of the image classification showing the behaviour of saline lakes over seasonal between 2017 and 2023 in the region of Gulf of Hammamet (Sebkhet de Sidi el Hani, Sebkha de Moknine and Sebkhet El Rharra). 

Figure \ref{fig06} demonstrates the outputs of the image classification showing variations in land cover types and saline lakes over seasonal between 2017 and 2023 in the region of Gulf of Gabès (Sebkhet en Noual). The land cover types were identified on the images using information adopted from and modified using the existing classification scheme on land cover types in Tunisia \cite{plants10102031} and data from Food and Agriculture Organization (FAO) \cite{FAO}. The existing land cover types for the country-level scale were modified according to the extent of the images covering study area and included the following types: 1) forest; 2) grasslands and rangelands; 3) sand and desert; 4) cultivated agriculture plantations and arable lands (olive tree, orchards, vine, palm groves); 5) urban areas and artificial land; 6) wasteland; 7) wetlands and water; 8) bare land and soil; 9) mixed vegetation; 10) saline soil. 

\begin{figure}[H]
\hspace{50pt}\makebox[0.90\linewidth][r]{%
	\begin{subfigure}[b]{.40\textwidth}
		\centering
			\includegraphics[width=0.95\textwidth]{Fig_05a.jpg}
		\caption{28 February 2017}
	\end{subfigure}%
	\begin{subfigure}[b]{.40\textwidth}
		\centering
		\includegraphics[width=0.95\textwidth]{Fig_05b.jpg}
		\caption{1 April 2017}
	\end{subfigure}%
	\begin{subfigure}[b]{.40\textwidth}
		\centering
		\includegraphics[width=0.95\textwidth]{Fig_05c.jpg}
		\caption{6 July 2017}
	\end{subfigure}%
}\\
\vfill \vspace{1mm}
\hspace{50pt}\makebox[0.90\linewidth][r]{%
	\begin{subfigure}[b]{.40\textwidth}
		\centering
			\includegraphics[width=0.95\textwidth]{Fig_05d.jpg}
		\caption{21 February 2023}
	\end{subfigure}%
	\begin{subfigure}[b]{.40\textwidth}
		\centering
		\includegraphics[width=0.95\textwidth]{Fig_05e.jpg}
		\caption{10 April 2023}
	\end{subfigure}%
	\begin{subfigure}[b]{.40\textwidth}
		\centering
		\includegraphics[width=0.95\textwidth]{Fig_05f.jpg}
		\caption{15 July 2023}
	\end{subfigure}%
}
\caption{Classification results showing land cover classes based on the processed Landsat images showing northern study area (Gulf of Hammamet): Sebkhet de Sidi el Hani (Sousse Governorate), Sebkha de Moknine (Mahdia Governorate) and Sebkhet El Rharra (Sfax Governorate): \textbf{(a)}: 28 February 2017; \textbf{(b)}: 1 April 2017; \textbf{(c)}: 6 July 2017; \textbf{(d)}: 21 February 2023; \textbf{(e)}: 10 April 2023; \textbf{(f)}: 15 July 2023.}\label{fig05}
\end{figure}

Changed seasons enabled to monitor and interpret how the salinisation process activate the accumulation of minerals in the salt lakes with changed colour of sabkhas. Thus, during the summer period, the structure of the crystals in the soil sediments in sabkhas was solidified by the increased evaporation as a result of the increased temperatures. In turn, the seasonal effects of increased temperature and lack of precipitation in the southern region (Figure \ref{fig06}) are reflected in the increased areas of sand and declined vegetation, as shown in the maps. Changes in salt lakes were determined at different years during recent decade by integrating Landsat images taken in winter/spring/summer months for comparison of periods when water inflow into Saharan lakes is high and low, respectively. Note that the decline in vegetation areas is clear on the images of 2017 from April to July showing sabkhas Sebkhet de Sidi el Hani, Sebkha de Moknine  and Sebkhet El Rharra, Figure \ref{fig06}. 

Compared to the difference between the February and April months, the decline in vegetation and increase in salinisation is more prominent, because it indicates the increase in summer heat temperatures which was contrasting compared to the spring months. In Figure \ref{fig06}, it is seen that the results are quite different for visualised series of maps for various seasons, because salt lakes contain different amounts of water with regard to salt, i.e., the classified images revealed various water/salt ratio in the sabkhas along with the development of heat during summer months and effects from Sahara desert in the southern study region. The classification of pixel confidence levels with rejection probability values indicating the accuracy of classification are presented in Figure \ref{fig07} for the northern region of the Gulf of Hammamet, and in Figure \ref{fig08} for the southern region of the Gulf of Gabès, respectively.

\begin{figure}[H]
\hspace{50pt}\makebox[0.90\linewidth][r]{%
	\begin{subfigure}[b]{.40\textwidth}
		\centering
			\includegraphics[width=0.95\textwidth]{Fig_06a.jpg}
		\caption{28 February 2017}
	\end{subfigure}%
	\begin{subfigure}[b]{.40\textwidth}
		\centering
		\includegraphics[width=0.95\textwidth]{Fig_06b.jpg}
		\caption{1 April 2017}
	\end{subfigure}%
	\begin{subfigure}[b]{.40\textwidth}
		\centering
		\includegraphics[width=0.95\textwidth]{Fig_06c.jpg}
		\caption{6 July 2017}
	\end{subfigure}%
}\\
\vfill \vspace{1mm}
\hspace{50pt}\makebox[0.90\linewidth][r]{%
	\begin{subfigure}[b]{.40\textwidth}
		\centering
			\includegraphics[width=0.95\textwidth]{Fig_06d.jpg}
		\caption{21 February 2023}
	\end{subfigure}%
	\begin{subfigure}[b]{.40\textwidth}
		\centering
		\includegraphics[width=0.95\textwidth]{Fig_06e.jpg}
		\caption{10 April 2023}
	\end{subfigure}%
	\begin{subfigure}[b]{.40\textwidth}
		\centering
		\includegraphics[width=0.95\textwidth]{Fig_06f.jpg}
		\caption{15 July 2023}
	\end{subfigure}%
}
\caption{Classification results showing land cover classes based on the processed Landsat images showing southern study area (Gulf of Gabès), Sebkhet en Noual (Sfax Governorate): \textbf{(a)}: 28 February 2017; \textbf{(b)}: 1 April 2017; \textbf{(c)}: 6 July 2017; \textbf{(d)}: 21 February 2023; \textbf{(e)}: 10 April 2023; \textbf{(f)}: 15 July 2023.}\label{fig06}
\end{figure} 

Varied temperature in seasonal climate periods has effects on the behaviour of lakes and salinisation of sabkhas which varies for different lakes. This also serves as indirect indicator of groundwater inflow and impacts of coastal climate on the eastern sabkhas, for instance for the case of Sebkha de Moknine which is located near the Bekalta coastal town. This lake experiences the least difference in the salinisation and was filled with water during the 2017, while in 2023 it experienced salinisation only during summer. In contrast, Sebkhet El Rharra experiences clear fluctuations in the salinisation during each year being regularly filled with water in winter and covered by salt crust during July. Such changes were detected both for year 2017 and for year 2023.

The extends of agricultural areas are much more similar in Figure \ref{fig06} for the Gulf of Gabès area, indicating that the plantation areas in southern regions are rather similar for different landscape spots compared to the northern segment (Figure \ref{fig05}). This is also seen in Figure \ref{fig06} for 2017 and 2023 years on the month of April. The use of comparative analysis of several images showing dynamics of land cover types as descriptors to reveal the development of aridification and desertification in southern regions closer to Sahara in study areas demonstrated RS approach to be a straightforward yet efficacious strategy to extract contextual information regarding land cover properties of landscapes. Thus, although RS data cannot directly detect landscape dynamics since this is dependent on the comparison of land cover properties for several scenes taken in different periods, nevertheless, processing of RS data enables to approximate the extent of landscape patches from classified images and evaluate water content in salt lakes as indirect indicators of salinisation during summer heat using image processing with the GIS methods. 

\begin{figure}[H]
\hspace{50pt}\makebox[0.90\linewidth][r]{%
	\begin{subfigure}[b]{.40\textwidth}
		\centering
			\includegraphics[width=0.95\textwidth]{Fig_07a.jpg}
		\caption{28 February 2017}
	\end{subfigure}%
	\begin{subfigure}[b]{.40\textwidth}
		\centering
		\includegraphics[width=0.95\textwidth]{Fig_07b.jpg}
		\caption{1 April 2017}
	\end{subfigure}%
	\begin{subfigure}[b]{.40\textwidth}
		\centering
		\includegraphics[width=0.95\textwidth]{Fig_07c.jpg}
		\caption{6 July 2017}
	\end{subfigure}%
}\\
\vfill \vspace{1mm}
\hspace{50pt}\makebox[0.90\linewidth][r]{%
	\begin{subfigure}[b]{.40\textwidth}
		\centering
			\includegraphics[width=0.95\textwidth]{Fig_07d.jpg}
		\caption{21 February 2023}
	\end{subfigure}%
	\begin{subfigure}[b]{.40\textwidth}
		\centering
		\includegraphics[width=0.95\textwidth]{Fig_07e.jpg}
		\caption{10 April 2023}
	\end{subfigure}%
	\begin{subfigure}[b]{.40\textwidth}
		\centering
		\includegraphics[width=0.95\textwidth]{Fig_07f.jpg}
		\caption{15 July 2023}
	\end{subfigure}%
}
\caption{Classification of pixel confidence levels with rejection probability values for the northern study area (Gulf of Hammamet): Sebkhet de Sidi el Hani (Sousse Governorate), Sebkha de Moknine (Mahdia Governorate) and Sebkhet El Rharra (Sfax Governorate): \textbf{(a)}: 28 February 2017; \textbf{(b)}: 1 April 2017; \textbf{(c)}: 6 July 2017; \textbf{(d)}: 21 February 2023; \textbf{(e)}: 10 April 2023; \textbf{(f)}: 15 July 2023.}\label{fig07}
\end{figure}

A time series analysis of the satellite images demonstrated in this paper characterises gradual changes in lakes over calendar periods and provides a robust view in environmental changes. Here, the obtained results presented a time series of processed satellite images that showed that the increase in salt lakes and desertification of surrounding areas is more distinct in the southern segment which showed that the Sebkhet en Noual is prone to salinisation even during the winter months, while the northern sabkhas Sebkhet de Sidi el Hani , Sebkha de Moknine and Sebkhet El Rharra show water-filled basins during winter. Furthermore, it demonstrated that the regions of Sfax are more prone to aridification compared to the areas of Mahdia and Sousse. 

Changes in water level in salt lakes were detected using RS data. The presented results can serve for environmental monitoring of Tunisian landscapes and analysis of environmental consequences of climate effects in central and north Tunisia. Although the results of image processing that analyse the environmental landscape properties and changes in the land cover types have been described earlier in various works on Tunisia \cite{Touhami,Khelifa,Hammami,Bounouh} providing a general pattern of land cover change by various GIS, the results proposed here are obtained using recent Landsat OLI/TIRS images processed by the GRASS GIS scripts. For instance, earlier works use the Landsat MSS/4 and SPOT HRV/3 images \cite{Park}, Landsat and SRTM processed by ArcGIS and ENVI commercial software \cite{Souissi,Afrasinei} or a combination of ETM+ and OLI 8 sensors of the Landsat products \cite{Mezned}. 

The reaction of materials sand salinisation processes during high evaporation soil results in changed structure of pores of soil with significantly reduced water content in soil material during periodic occurrences of a deep brine layer on the bottom of the lakes. Such behaviour of soil as lacustrine sediment layer demonstrates the response to climate and meteorological processes so that the considerable fluctuations in water level altered the soil structure in lacustrine sediments which is reflected as increased salinity and decreased water level, as evaluated during winter/spring/summer periods on satellite images. Moreover, the degree in salinisation of sabkhas well correlates with the location of each lake with southern lakes almost completely covered by salt even during summer, while northern lakes demonstrate fluctuation. 

\begin{figure}[H]
\hspace{50pt}\makebox[0.90\linewidth][r]{%
	\begin{subfigure}[b]{.40\textwidth}
		\centering
			\includegraphics[width=0.95\textwidth]{Fig_08a.jpg}
		\caption{28 February 2017}
	\end{subfigure}%
	\begin{subfigure}[b]{.40\textwidth}
		\centering
		\includegraphics[width=0.95\textwidth]{Fig_08b.jpg}
		\caption{1 April 2017}
	\end{subfigure}%
	\begin{subfigure}[b]{.40\textwidth}
		\centering
		\includegraphics[width=0.95\textwidth]{Fig_08c.jpg}
		\caption{6 July 2017}
	\end{subfigure}%
}\\
\vfill \vspace{1mm}
\hspace{50pt}\makebox[0.90\linewidth][r]{%
	\begin{subfigure}[b]{.40\textwidth}
		\centering
			\includegraphics[width=0.95\textwidth]{Fig_08d.jpg}
		\caption{21 February 2023}
	\end{subfigure}%
	\begin{subfigure}[b]{.40\textwidth}
		\centering
		\includegraphics[width=0.95\textwidth]{Fig_08e.jpg}
		\caption{10 April 2023}
	\end{subfigure}%
	\begin{subfigure}[b]{.40\textwidth}
		\centering
		\includegraphics[width=0.95\textwidth]{Fig_08f.jpg}
		\caption{15 July 2023}
	\end{subfigure}%
}
\caption{Classification of pixel confidence levels with rejection probability values for the southern study area (Gulf of Gabès): Sebkhet en Noual (Sfax Governorate: \textbf{(a)}: 28 February 2017; \textbf{(b)}: 1 April 2017; \textbf{(c)}: 6 July 2017; \textbf{(d)}: 21 February 2023; \textbf{(e)}: 10 April 2023; \textbf{(f)}: 15 July 2023.}\label{fig08}
\end{figure}

Similar to existing studies \cite{BAYATI2021126032,SHENG2016129,SCHRODER20221710}, the salinisation of lacustrine soil evaluated using RS data, this study showed the highest level of changes in the July month, and the most contrasting behaviour of brine accumulation in the lake of Sebkhet de Sidi el Hani (Sousse region) which has two connected sub-basins with different morphology hat might have effects on salinisation of the lake. The advanced scripting cartographic approach of GRASS GIS was introduced to evaluate the lacustrine properties of soil around the sabkhas in north Tunisia using RS data processing. It is well known that proprietary GIS and RS software are expensive since their development is costly (e.g., ArcGIS or Erdas Imagine). Moreover, very high resolution satellite images are not always easily accessible since they are provided on the commercial basis (e.g., SPOT, IKONOS). 

Therefore, a special advantage of this work is the use of the open source Landsat OLI/TIRS data and a free GRASS GIS software. This makes the present work repeatable, which can be continued in similar relevant studies in the future projects. The presented open source solutions on the optimal selection of data and tools can be supportive for environmental monitoring  of Tunisia. Hence, a series of the satellite images were processed with various combination of GRASS GIS modules and compared with the images covering identical areas on earlier periods, which was experimentally used for environmental monitoring of sabkhas in Tunisia. The performed experiments on RS data analysis supported monitoring of fluctuations in water level of salt lakes of Tunisia through the analysed time series analysis of satellite images. 

%----------- section ----------->
\section{Conclusions}

\subsection{Summary}

GRASS GIS software has been successfully applied to identify different stages in the landscape change in sabkhas of north Tunisia (with a special focus on several lakes -- Sebkhet de Sidi el Hani, Sebkha de Moknine, Sebkhet El Rharra and Sebkhet en Noual). The study also evaluated the effects of climate on seasonal changes in the lacustrine landscapes, as indicated by development of land cover types over the period of 2013 to 2023. Besides, the experiments served as a means of evaluating the climate effects as accelerators of salinisation development in sediments of lakes. The presented results confirm that GRASS GIS can be effectively used to processing and classification of RS imagery using programming approach that ensures automation and precision and accuracy of data processing. 

\subsection{Recommendations}

For future similar works on RS data analysis of salt lakes, we provide below the following recommendations and technical suggestions for mapping and image analysis:

\begin{enumerate}
	\item Make the initial study similar to the image processing workflow in this report, but also use other GRASS GIS modules applied for processing the target dates. For example, a useful approach is segmentation by 'i.segment'.
	\item Due to the  dynamics of plant phenology in different months, the period of image capture should be identical to evaluate changes for each image. For example, a correct comparison can be performed using images taken for a specific month (e.g., January) in a sequence of several years, or several identical months for two years, as presented in this study.
	\item Cloudiness vary a lot in different images. It is necessary to use a cloud-free series of images as an input dataset. While in this study, 10\% of cloud coverage was applied, the modified parameters optimised with different cloud content can be tested and applied with decreased cloud coverage. In this way, the use of 'i.landsat.acca' module of the GRSAS GIS enables Automatic Cloud Cover Assessment (ACCA) for a quick evaluation of image quality and suitability.
	\item Process other dataset in addition to the Landsat, e.g., Advanced Spaceborne Thermal Emission and Reflection Radiometer (ASTER) images. ASTER data can be proprocced using special modules of 'i.aster.toar', similar to the 'i.landsat.toar'. The main advantage of ASTER is a higher resolution--a pixel size of 15 m compared to the 30-m Landsat. 
	\item For mapping, take satellite images from the USGS on diverse dates to evaluate spatio-temporal dynamics and perform mapping on images captured on different seasonal periods to compare the difference in salinisation between the lakes.
	\item The  commercial GIS currently available on the market typically have license. Therefore, one can take open source GIS software such as QGIS or GRASS GIS, which have effective tools and extended functionality which are used for for RS and cartographic workload.
	\item Evaluate the integral effects from various factors that affect the landcover changes, i.e., the heat from the increased temperature, decreased precipitation and high evaporation lead to desertification of landscapes. 
	\item The access to the groundwaters and geographic location of the landscapes affect the behaviour of the vegetation. For instance, closeness to the desert or coastal areas affect differently vegetation patterns of the landscapes. 
	\item For analysis of vegetation health and vigour, the computation of the NDVI or other vegetation indices is a useful means of environmental assessment. In GRASS GIS, vegetation indices can be automatically calculated  by the software using the embedded formulae adjusted to the bands of the relevant satellite images. For example, NIR and Red channels correspond to the bands 5 and 4 for the Landsat 8-9 OLI/TIRS, while bands 4 and 3 for older Landsat 7 ETM+ images.
	\item Changes in land cover types within certain limits indicate that there is dynamics in landscapes as a response to the environmental and climate effects and processes. Moreover, it may indicate the cumulative effects of several factors on soil-vegetation conditions, including both climate-environmental processes and anthropogenic activities.
\end{enumerate}

\subsection{Concluding remarks}

This paper reports the results obtained from experimental investigations on satellite images Landsat 8-9 OLI/TIRS collected from the USGS. The tests on image processing were designed to evaluate the changes in several saline lakes of north Tunisia in various seasons in 2013 to 2023 and variations in water extent as indicators of salinisation. Specifically, we evaluated the proposed framework on RS data analysis using GRASS GIS software and a series of the satellite images from the USGS repository. We also proposed future directions in similar works to develop robust scheme on RS data processing tests in similar studies where cartographic approach to image analysis is applied in environmental monitoring projects aimed at evaluation of climate effects on landscape changes. In this work we also discussed the ways in which the scripting approach of cartographic data processing contrasts with traditional GIS methods of mapping that use Graphical User Interface (GUI), and complement to RS data processing for evaluating the connections between the climate effects and salinisation of sabkhas in Tunisia. 

We demonstrated the advantages of the GRASS GIS as a quantitative, robust and advanced cartographic tool for evaluating environmental properties based on image analysis. A workflow has been developed that uses a series of high-resolution Landsat 8-9 OLI/TIRS images with various combinations of GRASS GIS modules, evaluated to analyse the properties of lacustrine landscapes  of north Tunisia as a function of summer heat in Sahara and increase in evaporation during the period of 10 years (2013-2023). The approach of GRASS GIS demonstrated to be a more advanced way of image processing compared to traditional GIS software due to scripts that enables automation of workflow repeatability of the process. Thus, as showed in this study, the use of several module of GRASS GIS ensures multifold manipulations with images, such as in 'i.segment', 'i.maxlik' and other sub-tests that where applied for various steps of multispectral image analysis. 

A scripting approach of the GRASS GIS enabled us to map and analyse the rate of salinisation in sabkhas of Tunisia in real time in various seasons from 2013 to 2023 and to obtain the information from the processed images representing the land cover types in the surroundings. In turn, RS datasets enabled to model the behavior of landscapes in spring, summer and winter periods in north Tunisia which was used to compare the climate effects from rise in temperature on the evaporation of water in salt lakes, as directly related to heat increase during summer period in Sahara. The results demonstrated variations in water level in several sabkha and show technical effectiveness of the GRASS GIS in cartographic processing of the RS data for environmental tasks. Correlations between the increase in heat in various years and high salinisation of sabkhas were detected. The experiments on image processing reported in this article have investigated the salinisation behaviour in saline lakes with thermal variations using analysis of Landsat bands in multispectral channels. 

Due to the effectiveness and applicability of the performed experiments on image analysis by GRASS GIS -- flexibility of scripting approach and a wide variety of modules tuned for diverse tasks of image processing using its modular approach -- we expect the proposed method to be able to deal with other satellite imagery that cover diverse regions of north Africa to study development of lacustrine landscapes with effects from climate and fluctuations in temperature on Saharan ecosystems. Various sabkhas of north Tunisia with diverse surrounding soil types and salinisation/mineralisation degree were evaluated to generalise the proposed workflow of salt lake monitoring to a higher dimensional workload in the larger lakes and different world ecosystems with similar salt lakes. 

As demonstrated in this study, the RS data from the USGS EarthExplorer repository can be used to analyse the variations in salt lakes. It enables to derive information using data on spectral reflectance and brightness of pixels which strongly correlates with the distribution of various land cover types. In such a way, remote sensing data processing enables to detect patches in complex diversified lacustrine landscapes of north Sahara characterised with heterogeneous environment in contrasting climate of desert. This paper contributes to the environmental monitoring of Sahara, analysis of salt lake variations and regional studies of Tunisia, north Africa.

%----------- section ----------->

\funding{The publication was funded by the Editorial Office of \emph{Land}, Multidisciplinary Digital Publishing Institute (MDPI), by providing 100\% discount for the APC of this manuscript.}

\institutionalreview{Not applicable.}

\informedconsent{Not applicable.}

\dataavailability{Not applicable} 

\acknowledgments{The author thanks the reviewers for reading and review of this manuscript.}

\conflictsofinterest{The author declares no conflict of interest.} 

\abbreviations{Abbreviations}{
The following abbreviations are used in this manuscript:\\
\noindent 
\begin{tabular}{@{}ll}
ACCA & Automatic Cloud Cover Assessment \\
ASTER & Advanced Spaceborne Thermal Emission and Reflection Radiometer \\
FAO & Food and Agriculture Organization \\
GDAL & Geospatial Data Abstraction Library \\
GMT & Generic Mapping Tools \\
GIS & Geographic Information System \\
GUI & Graphical User Interface \\
GRASS & Geographic Resources Analysis Support System \\
Landsat OLI/TIRS & Landsat Operational Land Imager and Thermal Infrared Sensor \\
MODIS & Moderate Resolution Imaging Spectroradiometer \\ 
NIR & Near-infrared \\
RS & Remote Sensing \\
SPOT & Satellite Pour l’Observation de la Terre \\
SWIR & Short-wave infrared (SWIR)-1 \\
USGS & United States Geological Survey \\
\end{tabular}
}

%%%%%%%%%%%%%%%%%%%%%%%%%%%%%%%%%%%%%%%%%%
\begin{adjustwidth}{-\extralength}{0cm}
%\printendnotes[custom] % Un-comment to print a list of endnotes

\reftitle{References}

%=====================================
% References, variant A: external bibliography
%=====================================
\bibliography{Tunisia.bib}
%\nocite{*}

%%%%%%%%%%%%%%%%%%%%%%%%%%%%%%%%%%%%%%%%%%
\end{adjustwidth}
\end{document}
